\documentclass[11pt,a4paper]{article}
\usepackage[T1]{fontenc}
\usepackage[utf8]{inputenc}
\usepackage{lmodern}
\usepackage{microtype}
\usepackage[margin=2.3cm]{geometry}
\usepackage{amsmath}
\usepackage{booktabs}
\usepackage{tabularx}
\usepackage{enumitem}
\usepackage{graphicx}
\usepackage{xcolor}
\usepackage{float}
\usepackage[hidelinks]{hyperref}
\usepackage{fancyhdr}

\setlist{itemsep=2pt,topsep=4pt}
\setlength{\parindent}{0pt}
\setlength{\parskip}{5pt}
\definecolor{warn}{RGB}{170,40,20}
\definecolor{boxbg}{RGB}{245,245,240}

\newcommand{\um}{\,\textmu m}
\newcommand{\key}[1]{\fbox{\footnotesize\texttt{#1}}}
\newcommand{\file}[1]{\texttt{#1}}
\newcommand{\warning}[1]{\par\noindent\colorbox{boxbg}{\parbox{\dimexpr\linewidth-2\fboxsep}{\textcolor{warn}{\textbf{Warning.}} #1}}\par}
\newcommand{\note}[1]{\par\noindent\colorbox{boxbg}{\parbox{\dimexpr\linewidth-2\fboxsep}{\textbf{Note.} #1}}\par}
\newcommand{\fig}[3]{\begin{figure}[H]\centering\IfFileExists{figures/#1}{\includegraphics[width=#2\linewidth]{figures/#1}}{\fbox{\parbox{0.6\linewidth}{\centering figure missing: figures/#1}}}\caption{#3}\end{figure}}

\pagestyle{fancy}
\fancyhf{}
\lhead{DMD maskless lithography}
\rhead{Calibration guide}
\cfoot{\thepage}

\title{\textbf{DMD Maskless Lithography Setup}\\[4pt]\large Calibration and Operating Guide}
\author{}
\date{Version 1 --- November 2026}

\begin{document}
\maketitle
\thispagestyle{fancy}

\begin{abstract}
\noindent This guide describes how to align, calibrate and operate the maskless lithography setup built from an
InFocus IN124 DLP projector and a Zeiss Stemi 2000-C stereo microscope, exposing SU-8 at 405\,nm. It covers the
illumination optics, the projector display, the monitoring camera, focus, scale, power and drift, the grey response
of the projector, the flat-field (greyscale) correction, the dose calibration and the daily exposure procedure.
Numbers quoted as ``reference values'' are the results of the calibration in November 2026 (alignment v4); they
must be re-measured after any change of the optics.
\end{abstract}

\tableofcontents
\newpage

%==========================================================================
\section{Safety}
\begin{itemize}
  \item The source is a 405\,nm LED (Thorlabs M405L3). Near-UV/violet light is a blue-light hazard for the retina
        and is barely visible, so the eye does not protect itself by blinking. Follow the local rules for optical
        radiation.
  \item Never look into the eyepieces with the LED at exposure power. For visual alignment use either the green
        filter (Thorlabs FBH550-40; optical density about 5.8 at 405\,nm, so practically no violet light passes
        while enough green light is left to see the pattern) or the LED at its lowest setting
        (about 0.01--0.03\,mW at the sample), and keep viewing times short.
  \item Everything shown on the projector is projected onto the wafer, including menus and status text. Keep the LED
        off (or the green filter in) whenever a menu is open or a pattern is shown that should not expose the resist.
  \item Handle SU-8 and its developer according to their safety data sheets.
\end{itemize}

%==========================================================================
\section{System overview}
\subsection{Hardware}
\begin{table}[H]\centering\small
\begin{tabularx}{\linewidth}{@{}lX@{}}
\toprule
Component & Details \\
\midrule
Pattern generator & InFocus IN124 DLP projector, projection lens and integrator rod removed.
                    DMD 1024\,$\times$\,768 mirrors, 10.8\um\ pitch. Driven by a PC as a second display. \\
Projection optics & Zeiss Stemi 2000-C Greenough stereo microscope. The DMD sits in the image plane of the camera port;
                    the microscope images it demagnified onto the wafer. The two channels are tilted about $6^\circ$
                    to the vertical. \\
Scale on the wafer & 2.2\um\ per DMD pixel at zoom 5$\times$ (magnification $10.8/2.2 \approx 4.9$);
                     field $2.25\times1.69$\,mm. At zoom 2.5$\times$ about 4.4\um/px (calibrate). \\
Light source & Thorlabs M405L3 (405\,nm) with a single collector lens ($f = 50$\,mm, \O25\,mm) in a
               K\"ohler-type arrangement; an 8\,mm aperture right after the lens. \\
Resist & SU-8 TF 6000.5 (about 0.5\um\ film). \\
Power meter & Thorlabs PM100A with S120C photodiode sensor, wavelength set to 405\,nm, logged with the
              Optical Parameter Monitor software (CSV, 0.1\,s interval). \\
Camera & Logitech C920 (or C930e) without lens at the eyepiece image plane, manual exposure. \\
\bottomrule
\end{tabularx}
\end{table}

\subsection{Software}
\begin{table}[H]\centering\small
\begin{tabularx}{\linewidth}{@{}lX@{}}
\toprule
File & Purpose \\
\midrule
\file{test\_patterns/viewer.html} & Calibration viewer: test patterns, calibration grid (02), white (05),
      focus targets, dose ladder (D), uniformity probe (Q), grey ramp (R), correction preview. \\
\file{test\_patterns/camera.html} & Camera page: manual controls, histogram, clipping display, focus meter, scale bar. \\
\file{maskwriter/designer.html} & Mask designer: draws masks on the DMD pixel grid, exports 1024\,$\times$\,768 PNG. \\
\file{maskwriter/expose.html} & Exposure program: load a mask, set the time, greyscale correction on/off, expose, log. \\
\file{matlab\_greyscale/*.m} & MATLAB analysis: grey ramp, uniformity map with 3rd-order fit and correction
      overlay, comparison of runs with and without correction. \\
\bottomrule
\end{tabularx}
\end{table}

\subsection{Reference values (alignment v4, 3--4 November 2026)}
\begin{table}[H]\centering\small
\begin{tabular}{@{}ll@{}}
\toprule
Quantity & Value \\
\midrule
Usable area of the DMD & $x$ 0--896, $y$ 0--768\,px ($1.97\times1.69$\,mm); dark corner at the top right \\
Full-white power at the sample & 0.7--0.9\,mW (varies with LED drift; measure before use) \\
Field non-uniformity, no correction & 84--124\,\% of the mean (sd 9\,\%) in the usable area \\
Projector grey response & full light from grey 160; 155 $\to$ 91\,\%, 150 $\to$ 81\,\%, 140 $\to$ 61\,\% \\
Greyscale correction & target 86\,\%, grey levels 146--160, 87\,\% of the light kept \\
Working dose (SU-8 TF 6000.5) & 155--170\,mJ/cm$^2$ (nearly full thickness, 3\,px gaps open) \\
Resolution & dense lines/gaps $\geq$ 3\,px (6.6\um); 1\,px lines never resolve \\
\bottomrule
\end{tabular}
\end{table}

%==========================================================================
\section{Illumination alignment}
\label{sec:illum}
The LED is not imaged onto the DMD. Instead the collector lens images the LED into the entrance pupil of the
microscope (K\"ohler illumination): every point of the LED illuminates the whole DMD, so the field is uniform and the
light that leaves the DMD enters the microscope pupil.

\subsection{Geometry}
For the microscope the pupil seen from the camera port lies about $d \approx 320 \pm 40$\,mm beyond the DMD (measured).
With the lens a distance $L$ in front of the DMD, the LED must sit at the object distance
\begin{equation}
  u = \left(\frac{1}{f} - \frac{1}{L + d}\right)^{-1}
\end{equation}
in front of the lens, so that its image falls in the pupil. With $f = 50$\,mm, $L = 100$\,mm and $d = 320$\,mm this
gives $u \approx 57$\,mm. The size of the LED image in the pupil divided by the pupil size is the partial coherence
$\sigma$.

The lens rim acts as the field stop. The fully lit radius on the DMD and the width of the soft edge (penumbra) are
approximately
\begin{equation}
  r_\text{lit} \approx R_\text{lens}\,\frac{d}{L + d} - \frac{L}{u}\,a, \qquad
  w_\text{penumbra} \approx 2\,\frac{L}{u}\,a,
\end{equation}
with $R_\text{lens}$ the free radius of the lens (or of the 8\,mm aperture right after it) and $a \approx 0.75$\,mm the
half size of the LED emitter. A smaller $L$ (lens closer to the DMD) gives a larger lit area and a sharper edge; a
larger free aperture gives more margin around the DMD.

\subsection{Procedure}
\begin{enumerate}
  \item Mount lens and LED on a common axis pointing at the centre of the DMD. Start with $L \approx 100$\,mm and
        $u$ from equation (1).
  \item Show full white (pattern 05) at low LED power, green filter in. Look at the field in the microscope or with
        the camera: the DMD outline must be filled with light, without a dark corner. Remains of the projector
        housing (the old integrator-rod mount) can clip a corner: centre the beam or move the lens closer.
  \item Move the LED along the axis until the image of the LED is centred in the microscope pupil (look into the
        eyepiece tube from far away, or put a paper screen at the pupil position). The output power at the sample is
        maximal and the field most uniform when the LED image fills the pupil centrally.
  \item Fine-adjust the lateral position of lens/LED while watching the white field with the camera
        (Section~\ref{sec:camera}). Any bright rim or edge running across the field means the edge of the beam
        lies on the DMD: shift the beam away from that side, or enlarge the lit area.
  \item Measure the uniformity with the power meter (Section~\ref{sec:probe}). Repeat from step 3 if needed.
\end{enumerate}
\note{Uniformity and output power trade off: a larger lit area is more uniform but puts less light on the DMD. After
the v4 alignment the full-white power dropped by about half while the dark corner shrank to a small triangle at the
top right.}

%==========================================================================
\section{Projector display: 1:1 pixel mapping}
Every image pixel must land on exactly one DMD mirror.
\begin{enumerate}
  \item Connect the projector as a second display at 1024\,$\times$\,768, display scaling 100\,\%.
  \item Open the viewer (or \file{expose.html}) in Chrome, drag the window to the projector display, press
        \key{Ctrl+0} (browser zoom 100\,\%) and \key{F} (fullscreen).
  \item The menu shows ``Window 1024\,$\times$\,768, devicePixelRatio 1 --- 1:1 confirmed'' in green. If it is red,
        fix the scaling before going on; a scaled image blurs every pattern.
\end{enumerate}
The wafer image is rotated by $180^\circ$ relative to the screen (observed in the developed test patterns): the
top-right corner of the screen is at the bottom left on the wafer as seen in the microscope photos.

%==========================================================================
\section{Monitoring camera}
\label{sec:camera}
The camera sits without lens at the image plane of one eyepiece. Because of the Greenough geometry it looks through
the other channel than the projection: the specular reflection from the wafer goes back into the projection channel,
so the camera sees only light scattered by the wafer and the resist. The camera is therefore good for focus, position
and alignment, but \emph{not} a reliable reference for uniformity (its map did not match the power meter, $r\approx0.2$).

\subsection{Settings}
All automatic functions must be off. On Linux use \file{v4l2-ctl} (the device is the C920/C930e, not the built-in
camera; find it with \file{v4l2-ctl --list-devices}); on Windows open the DirectShow dialog
(\file{cam\_dialog\_windows.bat}) or use the controls in \file{camera.html}.
\begin{table}[H]\centering\small
\begin{tabular}{@{}ll@{}}
\toprule
Control & Setting \\
\midrule
Auto exposure & manual (aperture priority off) \\
Exposure time & so that the brightest part of the field is 80--90 counts in the blue channel \\
Gain & 0 \\
Dynamic frame rate & off (otherwise the exposure is changed behind your back) \\
Backlight compensation & 0 \\
White balance & manual, 4000\,K \\
Sharpness & 0 \\
Focus & manual (no lens: irrelevant, but keep autofocus off) \\
Power-line frequency & 50\,Hz \\
Zoom & 100 (no digital zoom) \\
LED & off \\
Brightness (black offset) & so that dark areas read 5--10 counts, never 0 \\
\bottomrule
\end{tabular}
\end{table}
\begin{itemize}
  \item The camera response is roughly linear only up to about 90 counts; above it the values are compressed and the
        field looks more uniform than it is. Pixels at 0 mean the black level is clipped, which exaggerates differences.
  \item The exposure time cannot exceed one frame period (33\,ms at 30\,fps); longer settings change nothing.
  \item Scale: about 0.81 camera pixel per DMD pixel (2.7\um\ per camera pixel), so the camera cannot resolve single
        DMD pixels.
\end{itemize}

%==========================================================================
\section{Focus}
\label{sec:focus}
\subsection{Depth of focus and the Greenough tilt}
The numerical aperture at the wafer is about 0.035--0.07 (port NA 0.007--0.015 times the magnification 4.9), which
gives a depth of focus of very roughly $\pm$40--80\um. Because the projection channel is tilted by $6^\circ$, the image
plane is tilted relative to a horizontal wafer: across the 2.25\,mm field the focus changes by about $\pm$120\um. The
centre can be in focus while the left and right edges are not; the developed test patterns indeed show softer edges on
one side. Re-focusing by $\Delta z$ also moves the pattern sideways by about $0.1\,\Delta z$ ($\tan 6^\circ$).
\begin{itemize}
  \item Put critical structures near the centre of the field.
  \item Lines whose spacing lies along the tilt direction blur first (in the dose ladder the horizontal line groups
        close before the vertical ones).
\end{itemize}

\subsection{Procedure}
There is no $z$ read-out, so focus is judged on the image.
\begin{enumerate}
  \item Green filter in (or LED at minimum). Show the focus target (\key{G} in the viewer) or the mask itself in
        alignment view (\key{O} in \file{expose.html}).
  \item Focus the wafer until the fine lines are sharpest, using the eyepiece or the focus meter of \file{camera.html}
        (gradient energy of the image; the peak and the 80\,\% line are marked).
  \item For best accuracy use a wafer that carries an old developed pattern and the split patterns \key{K}/\key{J}:
        one half of the screen is white and lights the old pattern, the other half projects fine lines. Focus the
        camera on the old pattern (the wafer surface), then adjust the wafer height until the projected lines are
        sharp in the camera, and repeat. The projected image is seen after reflection, so the error roughly halves
        with every round; two or three rounds are enough.
\end{enumerate}

%==========================================================================
\section{Scale calibration}
\label{sec:scale}
\begin{enumerate}
  \item Show the calibration grid (pattern 02: a cross every 64\,px) at low power and expose/develop it once, or
        view it directly.
  \item Measure the distance between crosses far apart (e.g. 640\,px) with a calibrated microscope reticle or the
        scale bar of an imaging microscope.
  \item Scale $s$ = measured distance / number of pixels. Nominal $s = 2.2\um$/px at 5$\times$. Enter $s$ in the
        designer (\emph{Scale}) and in \file{expose.html} (\emph{Scale on the wafer}).
\end{enumerate}
Repeat for every zoom setting you use; the zoom knob is not precise enough to rely on the nominal value.

%==========================================================================
\section{Power, irradiance and drift}
\label{sec:power}
\subsection{Power at the sample}
\begin{enumerate}
  \item Place the S120C sensor at the sample position, centred under the field; the 9.5\,mm aperture collects the
        whole field. Wavelength 405\,nm.
  \item Show full white: \key{W} in \file{expose.html} (with the correction set as you will expose) or pattern 05 in
        the viewer. Read the power $P$.
  \item Irradiance (average over the full field):
        \begin{equation}
          E = \frac{P}{A_\text{field}}, \qquad A_\text{field} = (1024\,s)(768\,s) = 0.0381\,\text{cm}^2 \text{ at } s = 2.2\um/\text{px}.
        \end{equation}
        Example: $P = 0.7$\,mW $\Rightarrow E = 18.4$\,mW/cm$^2$.
\end{enumerate}

\subsection{Drift}
The LED output changes with temperature: up to 10--20\,\% within a quarter of an hour was observed.
\begin{itemize}
  \item Let the LED run for 15--20\,min before measuring or exposing.
  \item Once, log full white for 30\,min from a cold start to learn how long the warm-up takes on this setup. If the
        power still moves after warm-up, check the mounts (thermal movement of the beam on the DMD/pupil).
  \item Measure the power immediately before each exposure and set the time from it (Section~\ref{sec:dose}).
  \item For uniformity runs use the drift reference tile (Section~\ref{sec:probe}); otherwise drift appears as a
        top--bottom gradient, because the rows are measured one after the other.
\end{itemize}

%==========================================================================
\section{Grey response of the projector}
\label{sec:ramp}
The correction dims bright regions with grey levels. A DLP projector makes grey by switching the mirrors on and off
within each frame, and its transfer curve is not a simple gamma. It has to be measured.

\subsection{Measurement}
\begin{enumerate}
  \item Power meter at the sample, logging to CSV. Corner text off (\key{H}).
  \item In the viewer set \emph{Grey ramp levels} (default: 255, 200, 180, 175, 170, 165, 160, 155, 150, 140, 130,
        120, 110, 100, 90, 80, 70, 60, 50), probe time about 3\,s and dark time about 1\,s.
  \item Arm the ramp (\key{R}) and press \key{Space}. A 256\,px square in the centre is shown at each level in turn.
  \item Analyse the log with \file{analyze\_grey\_ramp.m} (set \file{file} and \file{levels}). It prints the table,
        the knee and the string \file{level:power, \ldots} for the field \emph{Grey curve} in \file{expose.html} and the
        viewer, and saves the curve for the correction.
\end{enumerate}

\fig{grey_ramp.png}{0.95}{Grey response measured on 3 Nov 2026. Full light from grey 160 upwards; below it about
2\,\% light per grey level; a shelf between 100 and 140.}

\subsection{Reference result}
Full light from grey 160 (plateau, within the noise of 2\,\%), 155 $\to$ 91\,\%, 150 $\to$ 81\,\%, 140 $\to$ 61\,\%,
100--140 a shelf at 51--61\,\%, 90 $\to$ 37\,\%, 70 $\to$ 14\,\%, 50 $\to$ 3\,\%. Everything above 160 gives the
same light, which is why a correction that only used grey levels above 160 had no effect. The correction works in the
steep range 145--160.
\warning{Re-measure the ramp after any change of projector settings or firmware state; the curve depends on them.}

%==========================================================================
\section{Uniformity map}
\label{sec:probe}
\subsection{Measurement with the probe}
\begin{enumerate}
  \item LED warmed up, power meter at the sample, logging to CSV, corner text off (\key{H}), correction off.
  \item Viewer settings: \emph{Probe grid} $16\times12$, \emph{tiles cover 100\,\%} on, time per tile 3\,s, dark time
        about 0.8\,s, \emph{drift reference tile after every 16 tiles} (one per row).
  \item Arm the probe (\key{Q}), press \key{Space}. The tiles of $64\times64$\,px are shown one by one, row by row from
        the top left, with a short beep at each tile and a different beep for reference tiles. A run takes about
        14\,min.
  \item Optionally download the probe CSV from the viewer menu (times of every tile) for your records.
\end{enumerate}

\subsection{Analysis and correction (MATLAB)}
\begin{enumerate}
  \item Run \file{analyze\_grey\_ramp.m} first (Section~\ref{sec:ramp}).
  \item In \file{analyze\_probe.m} set \file{file}, \file{cols}/\file{rows}, \file{refEvery}, the usable area and the
        \file{target}. Run it.
  \item The script finds every tile in the log (tiles too dim to detect are taken at their predicted time), subtracts
        the dark level, divides out the drift, and fits a 2-D polynomial of order 3 to the tiles in the usable area:
        \begin{equation}
          I(X,Y) = \sum_{i+j\le3} a_{ij}\,X^iY^j, \qquad X = \frac{x - 448}{448},\; Y = \frac{y - 384}{384},
        \end{equation}
        i.e.\ 10 terms. Order 3 is chosen because its residual (2.1\,\%) equals the measurement noise; order 2 leaves
        3.8\,\%, order 4 gains little and starts to follow noise. Tiles outside the usable area get the value of the
        nearest usable tile.
  \item Outputs:
        \begin{itemize}
          \item \file{map\_16x12\_for\_viewer.txt}: the smooth map in \% of the mean;
          \item \file{correction\_overlay\_1024x768.png}: the grey level of every DMD pixel;
          \item plots of the measured map, the fit and the residual.
        \end{itemize}
  \item Paste the map into \emph{Calibration data $\to$ correction map} in \file{expose.html} (and the viewer) and set the
        same \emph{Target level}. Paste the grey curve string into \emph{Grey curve}.
\end{enumerate}

\fig{probe_maps.png}{0.95}{Probe without correction (3 Nov, 17:05), its 3rd-order fit, the residual and the expected
corrected field. Outlined tiles lie outside the usable area.}

\subsection{How the correction is applied}
For every DMD pixel the viewer interpolates the map bilinearly between the tile centres (value $v$, mean of the map
$\bar v$), computes the required dimming and looks up the grey level in the inverted grey curve $L(g)$:
\begin{equation}
  c = \min\!\left(1,\; \frac{\text{target}}{100\,v/\bar v}\right), \qquad g = L^{-1}(c).
\end{equation}
Lit pixels of a mask are shown with grey $g$ instead of 255; pixels already at or below the target stay at full light.
The light kept on average is about the target (87\,\% for target 86), so the exposure time for the same dose is
$\approx 1.15\times$ longer --- this is taken care of automatically if the power is measured with the correction on.

\subsection{Verification}
Run the probe again with the correction \emph{on} (same grid and times), then run \file{compare\_correction.m} with the
run without correction (\file{fileBefore}) and the new run (\file{fileAfter}). It shows, on one colour scale, the
deviation from the mean without correction, its 3rd-order fit, and the measured result with correction, and prints
range, standard deviation, within-row spread (free of drift) and tile noise. A good correction leaves only a few
percent, comparable to the noise.

%==========================================================================
\section{Dose calibration}
\label{sec:dose}
\subsection{Dose ladder}
\begin{enumerate}
  \item Spin and bake SU-8 TF 6000.5 as usual.
  \item Measure the power with the correction on as it will be used. Choose a base time and a step so that the eight
        cells span about 110--200\,mJ/cm$^2$ (e.g.\ at 0.7\,mW: base 6\,s, step 0.5\,s gives 6.5--10\,s).
  \item In the viewer: dose ladder (\key{D}), set step time and base time (\key{M}), \key{Space}. Cell $n$ receives
        base $+\,n\times$step.
  \item Post-exposure bake, develop, inspect. For each cell judge: film thickness (interference colour; full thickness
        is the same colour as large exposed areas), separation of the dense line groups (1--10\,px), isolated lines and
        dots.
\end{enumerate}

\fig{dose_ladder.jpg}{0.7}{Developed dose ladder (4 Nov 2026, base 6\,s, step 0.5\,s, 0.7\,mW with correction; image
rotated $180^\circ$ relative to the screen: cell 1$\times$ bottom right, 8$\times$ top left).}

\subsection{Reference result (4 Nov 2026)}
\begin{table}[H]\centering\small
\begin{tabular}{@{}llll@{}}
\toprule
Cells & Time & Dose & Result \\
\midrule
1$\times$--2$\times$ & 6.5--7.0\,s & 120--130\,mJ/cm$^2$ & thin film; 3\,px clean, 2\,px nearly resolved \\
3$\times$--4$\times$ & 7.5--8.0\,s & 139--148\,mJ/cm$^2$ & below full thickness; 3\,px clearly resolved \\
5$\times$--6$\times$ & 8.5--9.0\,s & 157--167\,mJ/cm$^2$ & nearly full thickness; 3\,px gaps still open \\
7$\times$--8$\times$ & 9.5--10\,s & 176--185\,mJ/cm$^2$ & full thickness; narrow gaps close \\
\bottomrule
\end{tabular}
\end{table}
Working dose: \textbf{155--170\,mJ/cm$^2$}. With the power $P$ measured with the correction on:
\begin{equation}
  t = \frac{D}{P/A_\text{field}} \qquad (0.7\text{\,mW}, 165\text{\,mJ/cm}^2 \Rightarrow 9.0\text{\,s}).
\end{equation}
\note{One ladder step of 0.5\,s is 5--8\,\% dose, less than the field non-uniformity without correction; the cell
boundaries in the table are uncertain by about one step.}

\subsection{Resolution and design rules}
\fig{siemens_star.jpg}{0.55}{Developed Siemens star (36 spokes). The dark gaps close below about 10\,px design width
near large exposed areas.}
\begin{itemize}
  \item Dense lines and gaps: $\geq 3$\,px (6.6\um). 1\,px lines never resolve.
  \item Isolated lit lines: 3--4\,px to print reliably; isolated 2\,px dots print at the working dose but round off.
  \item Exposed areas grow into neighbouring dark areas by about 2--3\,px per edge, most next to large exposed areas
        (Siemens star: gaps narrower than about 10\,px closed, the central hole shrank by 3\,px in radius). Make dark
        gaps between large lit areas wider in the mask, or reduce the dose.
  \item Keep critical structures in the centre of the field (focus tilt, Section~\ref{sec:focus}) and inside the usable area.
\end{itemize}

%==========================================================================
\section{Making masks}
\begin{enumerate}
  \item Open \file{maskwriter/designer.html}. Set the scale (Section~\ref{sec:scale}) and check the usable area.
  \item Draw with rectangles, circles, lines (width), polygons and text. Every object either adds light or cuts it away.
        Dimensions are entered in \um\ and rounded to whole pixels; the design grid with snap helps with regular
        layouts.
  \item Watch the checks panel: lit features under 3\,px and dark gaps under 4\,px are flagged.
  \item Save the design (\file{.mask.json}), export the PNG (1024\,$\times$\,768, white = exposed), or send it
        directly to the exposure window.
\end{enumerate}

%==========================================================================
\section{Daily exposure procedure}
\begin{enumerate}
  \item Switch on the LED (green filter in) and let it warm up 15--20\,min.
  \item Open \file{expose.html} on the projector display; check the 1:1 message; \key{F} for fullscreen.
  \item Check that the correction map, target and grey curve are the current ones (Calibration data).
  \item Power meter at the sample, full white with \key{W}, correction as it will be used; enter $P$ in the dose helper;
        set the time from the dose (155--170\,mJ/cm$^2$).
  \item Load the mask (file or from the designer). Check the preview: inside the usable area, correct orientation.
  \item Place the wafer, focus with the alignment view \key{O} (green filter / minimum LED), then return to the black
        armed screen.
  \item LED to exposure power, filter out. \key{Space} exposes; \key{Esc} aborts immediately.
  \item Copy the log (time, measured exposure time, correction, dose) into the lab book.
\end{enumerate}

%==========================================================================
\section{Troubleshooting}
\begin{table}[H]\centering\small
\begin{tabularx}{\linewidth}{@{}>{\raggedright\arraybackslash}p{0.33\linewidth}X@{}}
\toprule
Symptom & Likely cause and remedy \\
\midrule
Pattern blurred everywhere & Window not 1:1 (scaling, zoom); wafer out of focus. \\
One side of the field softer & Greenough tilt: focus on the centre, keep critical parts central. \\
Dark corner / one edge dimmer & Beam edge on the DMD: re-centre the illumination (Section~\ref{sec:illum}); keep the
                                usable area. \\
Correction has no effect & Grey levels above the knee (160) give full light; re-measure the grey ramp and update the
                           curve. \\
Corrected field still has a top--bottom gradient & LED drift during the probe; use reference tiles and a warm LED. \\
Narrow gaps closed, features too wide & Dose too high or flare from large lit areas: lower dose, wider gaps. \\
Thin, coloured, partly removed film & Under-exposed: check the power and the time; LED drift. \\
Camera image looks flat or extreme & Camera saturated ($>$90 counts) or black clipped (0 counts); fix exposure and
                                     brightness. \\
Power-meter log cannot be analysed & Corner text left on, steps missing or extra: switch the text off (\key{H}), check
                                     the number of steps and the settings in the script. \\
\bottomrule
\end{tabularx}
\end{table}

%==========================================================================
\appendix
\section{Calibration record}
\begin{table}[H]\centering\small
\begin{tabular}{@{}p{0.42\linewidth}p{0.25\linewidth}p{0.25\linewidth}@{}}
\toprule
Item & Value & Date / initials \\
\midrule
Lens--DMD distance $L$, LED distance $u$ & & \\
Scale at 5$\times$ ($\mu$m/px) & & \\
Usable area ($x_0$--$x_1$, $y_0$--$y_1$) & & \\
Full-white power, correction off / on (mW) & & \\
LED warm-up time to $\pm$2\,\% (min) & & \\
Grey ramp file, knee & & \\
Probe file without correction, sd (\%) & & \\
Correction target (\%), grey range & & \\
Probe file with correction, sd (\%) & & \\
Dose ladder: working dose (mJ/cm$^2$) & & \\
Camera exposure / brightness & & \\
\bottomrule
\end{tabular}
\end{table}

\section{Keys}
\begin{table}[H]\centering\small
\begin{tabularx}{\linewidth}{@{}lX@{}}
\toprule
\file{expose.html} (armed screen) & \key{Space} expose, \key{Esc} abort/menu, arrows/PgUp/PgDn/\key{T} time,
  \key{X} correction, \key{O} alignment view, \key{W} full white, \key{V} preview, \key{H} text, \key{S} sound,
  \key{F} fullscreen \\
Viewer (armed screen) & pattern keys 1--9, A, B, C, L, G, K, J, D, Q, R, U; \key{M} base time of the dose ladder;
  \key{X} correction; \key{O} focus mode \\
\file{designer.html} & \key{V} select, \key{R} rectangle, \key{C} circle, \key{L} line, \key{P} polygon, \key{T} text,
  \key{H} pan, \key{G} pixel grid, \key{Shift+G} design grid, \key{Ctrl+Z}/\key{Ctrl+Y}, \key{Ctrl+D}, \key{Del},
  arrows nudge \\
\bottomrule
\end{tabularx}
\end{table}

\end{document}
